For $A\subseteq[1,N]$ let $S_r(A)$ be its sums of $r$ distinct elements. Call $A$ admissible if $S_r(A)\cap S_s(A)=\varnothing$ for $r\ne s$. We prove that the maximal size $k(N)$ satisfies $k(N)\sim2\sqrt N$. Our explicit external input is the structural Theorem 2 quoted by Deshouillers and Freiman in \emph{On an additive problem of Erd\H{o}s and Straus, 2}, Ast\'erisque 258 (1999), \url{https://archive.numdam.org/item/AST_1999__258__141_0.pdf}. We do not assume their final extremal bound.

The structural input says that if an admissible $A\subseteq[1,N]$ has $|A|>1.96\sqrt N$ and $N$ is sufficiently large, there are $C\subseteq A$, an integer $q\ge1$, and a fixed-cardinality subset-sum set of $C$ containing an arithmetic progression $J$ of at least $3N^{5/6}$ terms with difference $q$, such that
\[
 |C|\le10^5N^{5/12},\qquad B=A\setminus C
 \text{ lies in a progression of difference }q\text{ with at most }N^{7/12}\text{ terms}.\tag{1}
\]
We derive the required sharp leading constant from this more structural statement.
